\chapter{Projection and Representation}
\index{projection}\index{representation}

\section{Why This Projection Rather Than Another?}

Parts~One through Four largely treated the compression operator $\pi$ as given. Chapter~15 asks the prior question: why is one compression chosen rather than another?

Different compressions of the same trajectory space create different worlds of recoverable history. Two observers compressing the same $\mathcal{X}$ using different operators inherit different equivalence relations, different query families for which reconstruction is possible, and different failure modes.

\section{Projection Creates Worlds}

The compression $\pi : \mathcal{X} \to \mathcal{M}$ partitions $\mathcal{X}$ into fibers $\pi^{-1}(m)$. The recovered world is not $\mathcal{X}$ but the quotient $\mathcal{X}/{\sim_\pi}$---the world of equivalence classes, where some differences have been erased and others preserved.

Memory preserves semantic content while discarding perceptual detail. Archives preserve formal reasoning while discarding social context. Physical dynamics preserve macroscopic observables while discarding microscopic states. Scientific theories preserve invariant structure while discarding measurement noise. In each case the projection \emph{creates} the world available to reconstruction.

\section{Projection as Structured Forgetting}

A representation is a pattern of preserved distinctions. Projections forget selectively: they collapse some distinctions while preserving others. What a representation says is determined not by what it contains but by what it \emph{excludes}---which trajectories it rules out.

Every domain in this book exemplifies this: sketches preserve derivative structure; branch trees preserve admissible branch structure; skeletons preserve accessible future structure; memory preserves semantic-narrative structure; archives preserve doctrinal structure.

\section{The Pseudoinverse Perspective}

When $\pi$ is many-to-one, exact inversion is impossible. The query-relative reconstruction operator $\pi^+_Q : \mathcal{M} \to \mathbb{R}$ is a \emph{query-relative pseudoinverse} \cite{Penrose1955,BenIsrael2003}. It selects a canonical representative from each fiber and evaluates the query on it.

Reconstruction is always about the projection of the past, never the past directly. The accuracy of reconstruction depends not on choosing the correct representative but on whether all representatives in the fiber produce approximately the same query answer---which is precisely the stability condition.

\section{Alignment as Geometric Property}

An aligned projection collapses trajectories in \emph{stable directions} (where $\sigma_r$ decays rapidly) and preserves trajectories in \emph{important directions} (where $\sigma_r$ remains large at low $r$). Geometrically: the projection's null space is parallel to the stable directions.

Misalignment is geometrically transparent: sufficiency failure occurs when the projection collapses important directions; over-compression when it preserves stable directions unnecessarily.

\section{Representation as Constraint}

Representations are not primarily depictions. They are constraints: specifications of which trajectories remain possible. The significance of a representation is what it excludes, not what it contains. This constraint-based view explains why representations can be accurate without being complete, why multiple representations can be equally valid for the same query family, and why representation precedes inquiry.

\section{Projection Before Knowledge}

Before a system can know anything, it must project. Memory must compress before it can retrieve. Archives must be built before they can be consulted. Theories must model before they can predict. The world of recoverable inquiry is always defined before the inquiry begins. Knowledge is reconstruction performed inside a projection-induced quotient space.

\textbf{Stability determines when compression is possible. Projection determines what is preserved. Together they define the geometry of recoverable pasts.}
